%! TEX root = part1.tex
% the main TeX file which is intended to compile, :VimtexReload after adjustment
   % this file should be included in the main file
% :h vimtex-tex-root
\documentclass[../skb.tex]{subfiles}
% ensure the class of docs
% enumerate may not work

\begin{document}

\section{Pendahuluan}%
\label{sec:Pendahuluan}


\begin{frame}[t,mybg=pg_1,mytitle=standard,dark]
\end{frame}

\begin{frame}[t,mybg=pg_2,mytitle=standard,dark]
\end{frame}

\begin{frame}[t,mycolor=digiPH_leaf,mytitle=standard,dark]
	\frametitle{Pendahuluan}

	Bangunan tinggi (high-rise building) memerlukan perencanaan yang kompleks karena:

	\begin{itemize}
		\item Skala vertikal yang besar
		\item Beban struktur yang signifikan
		\item Risiko keselamatan yang tinggi
		\item Sistem utilitas yang kompleks
	\end{itemize}

\end{frame}
%---------------------------------

\begin{frame}[t,mycolor=digiPH_leaf,mytitle=standard,dark]
	\frametitle{1. Fungsi Struktural}

	Bangunan harus mampu:

	\begin{itemize}
		\item Menahan beban vertikal (beban mati dan hidup)
		\item Menahan beban lateral (angin dan gempa)
		\item Menjamin stabilitas dan kekakuan struktur
	\end{itemize}

	Sistem struktur umum:
	\begin{itemize}
		\item Rigid Frame
		\item Shear Wall
		\item Core System
		\item Tube System
		\item Outrigger System
	\end{itemize}

\end{frame}
%---------------------------------

\begin{frame}[t,mycolor=digiPH_leaf,mytitle=standard,dark]
	\frametitle{2. Fungsi Keselamatan}

	\begin{itemize}
		\item Sistem proteksi kebakaran
		\item Jalur evakuasi dan tangga darurat
		\item Sistem tahan gempa
		\item Sistem penangkal petir
		\item Kontrol asap
	\end{itemize}

\end{frame}
%---------------------------------

\begin{frame}[t,mycolor=digiPH_leaf,mytitle=standard,dark]
	\frametitle{3. Fungsi Sirkulasi}

	Sirkulasi vertikal dan horizontal:

	\begin{itemize}
		\item Lift penumpang dan servis
		\item Tangga darurat
		\item Koridor dan akses publik
		\item Zonasi lift (low, mid, high zone)
	\end{itemize}

\end{frame}
%---------------------------------

\begin{frame}[t,mycolor=digiPH_leaf,mytitle=standard,dark]
	\frametitle{4. Fungsi Utilitas dan Mekanikal}

	\begin{itemize}
		\item Sistem listrik
		\item Air bersih dan air kotor
		\item HVAC
		\item Sistem komunikasi
		\item Pengelolaan limbah
	\end{itemize}

	Biasanya terdapat \textit{mechanical floor}.

\end{frame}
%---------------------------------

\begin{frame}[t,mycolor=digiPH_leaf,mytitle=standard,dark]
	\frametitle{5. Fungsi Kenyamanan}

	\begin{itemize}
		\item Kenyamanan termal
		\item Pencahayaan alami
		\item Pengendalian kebisingan
		\item Kontrol getaran akibat angin/gempa
	\end{itemize}

\end{frame}
%---------------------------------

\begin{frame}[t,mycolor=digiPH_leaf,mytitle=standard,dark]
	\frametitle{6. Fungsi Efisiensi dan Ekonomi}

	\begin{itemize}
		\item Efisiensi luas lantai (net to gross ratio)
		\item Efisiensi struktur dan material
		\item Efisiensi energi (green building concept)
		\item Biaya konstruksi dan operasional
	\end{itemize}

\end{frame}
%---------------------------------

\begin{frame}[t,mycolor=digiPH_leaf,mytitle=standard,dark]
	\frametitle{7. Fungsi Estetika dan Identitas}

	\begin{itemize}
		\item Bentuk arsitektural ikonik
		\item Integrasi dengan konteks kota
		\item Identitas visual
		\item Konsep desain yang kuat
	\end{itemize}

\end{frame}
%---------------------------------

\begin{frame}[t,mycolor=digiPH_leaf,mytitle=standard,dark]
	\frametitle{Kesimpulan}

	Perancangan bangunan tinggi harus mempertimbangkan:

	\begin{itemize}
		\item Struktur
		\item Keselamatan
		\item Sirkulasi
		\item Utilitas
		\item Kenyamanan
		\item Efisiensi
		\item Estetika
	\end{itemize}

	Semua aspek saling terintegrasi.

\end{frame}
%---------------------------------


\begin{frame}[c,mybg=pg_3,mytitle=standard,dark]\end{frame}
\begin{frame}[c,mybg=pg_4,mytitle=standard,dark]\end{frame}

\begin{frame}[t,mycolor=digiPH_leaf,mytitle=standard,dark]
	\frametitle{Parameter Pemilihan Sistem Struktur}

	\begin{itemize}

		\item Pertimbangan ekonomi: construction cost, nilai kapitalisasi,
		      rentable space variation, dan cost of time variation.

		\item Kecepatan konstruksi: profil bangunan, pengalaman dan expertise,
		      metode pelaksanaan, material struktur, tipe konstruksi (cast-in-situ,
		      precast, kombinasi), serta kondisi industri konstruksi lokal.

		\item Geometri bangunan: panjang, lebar, tinggi, serta vertical profile
		      dan bentuk massa bangunan.

		\item Pembatasan ketinggian (height restriction) dan tingkat kelangsingan
		      (slenderness ratio).

		\item Konfigurasi denah: depth–width ratio dan tingkat regularitas
		      (UBC/NEHRP).

		\item Kinerja struktur: kekuatan, kekakuan (lentur, geser, torsi),
		      dan daktilitas (strain, curvature, displacement).

		\item Jenis pembebanan: gravitasi, angin, gempa, dan beban khusus lainnya.

		\item Kondisi tanah pendukung bangunan.

	\end{itemize}

\end{frame}

\section{Aplikasi SK Bangunan Tinggi}%
\label{sec:Aplikasi SK Bangunan Tinggi}

\begin{frame}[c,mybg=pg_5,mytitle=standard,dark]\end{frame}


\begin{frame}[t,mycolor=digiPH_gray,mytitle=standard,light]
	\frametitle{Klasifikasi Sistem Struktur Bangunan Tinggi}

	\begin{itemize}

		\item Dinding pendukung sejajar (Parallel bearing wall)
		\item Inti dan dinding pendukung fasade (Core and facade bearing wall)
		\item Boks berdiri sendiri (Self-support box)
		\item Plat terkantilever (Cantilevered slab)
		\item Plat rata (Flat slab)
		\item Interspasial (Interspatial system)
		\item Rangka dan dinding geser (Shear wall – frame interaction system)
		\item Rangka selang seling (Staggered truss)
		\item Rangka kaku (Rigid frame)
		\item Rangka kaku dan inti (Rigid frame and core)
		\item Rangka trussed (Trussed frame)
		\item Rangka belt trussed dan inti (Belt trussed frame and core)
		\item Tabung dalam tabung (Tube in tube)

	\end{itemize}

\end{frame}

\section{Strategi Perancangan}%
\label{sec:Strategi Perancangan}

\begin{frame}[c,mybg=pg_6_1,mytitle=standard,dark] \end{frame}
\begin{frame}[c,mybg=pg_6_2,mytitle=standard,dark] \end{frame}
\begin{frame}[c,mybg=pg_6_3,mytitle=standard,dark] \end{frame}

\section{Prinsip Dasar SK Bangunan Tinggi}%
\label{sec:Prinsip Dasar SK Bangunan Tinggi}

\begin{frame}[c,mybg=pg_7_1,mytitle=standard,dark] \end{frame}
\begin{frame}[c,mybg=pg_7_2,mytitle=standard,dark] \end{frame}
\begin{frame}[c,mybg=pg_7_3,mytitle=standard,dark] \end{frame}

\begin{frame}[label=current,c,mybg=pondasi,mytitle=standard,dark] \end{frame}

%================================================
\begin{frame}[t,mycolor=digiPH_leaf,mytitle=standard,dark]
	\frametitle{Sistem Struktur Berbasis Dinding}

	\begin{itemize}
		\item Dinding pendukung sejajar (Parallel bearing wall)
		\item Core dan dinding fasade (Core and façade bearing wall)
		\item Boks berdiri sendiri (Self-support box)
	\end{itemize}

\end{frame}
%================================================


%================================================
\begin{frame}[t,mycolor=digiPH_leaf,mytitle=standard,dark]
	\frametitle{Sistem Struktur Berbasis Plat}

	\begin{itemize}
		\item Plat terkantilever (Cantilevered slab)
		\item Plat rata (Flat slab)
		\item Interspasial (Interspatial system)
	\end{itemize}

\end{frame}
%================================================


%================================================
\begin{frame}[t,mycolor=digiPH_leaf,mytitle=standard,dark]
	\frametitle{Sistem Rangka dan Struktur Tinggi}

	\begin{itemize}
		\item Rangka selang-seling (Staggered truss)
		\item Gantung (Suspension system)
		\item Rangka kaku (Rigid frame)
		\item Rangka kaku dan inti (Rigid frame and core)
		\item Rangka trussed (Trussed frame)
		\item Rangka belt trussed dan inti (Belt trussed frame and core)
		\item Tabung dalam tabung (Tube in tube)
		\item Kumpulan tabung (Bundled tube)
	\end{itemize}

\end{frame}
%================================================

\begin{frame}[t,mybg=pg_8_1,mytitle=standard,dark]\end{frame}
\begin{frame}[c,mycolor=digiPH_white,mytitle=standard,light]
	\scalebox{.9}{
		\bgimg{1-}{pg_8_2}}
\end{frame}
\begin{frame}[t,mybg=pg_8_3,mytitle=standard,dark]\end{frame}
\begin{frame}[t,mybg=pg_8_4,mytitle=standard,dark] \end{frame}










\bibliographystyle{apalike}
\bibliography{../biblio.bib} % required for citecompletion, not work in mainfile
\end{document}

